
\texttt{LABEL\_NAME\_GLOB} ($31$) is the one clean expressiveness case left, and with
\texttt{.weighted()} withdrawn, the disjunctions reclassified and
\texttt{FILL\_INTERPOLATION} returned by Remark~\ref{rem:fill}, it and the percentiles
and the empty-valued scopes are all that remains. A wildcard in a tag \emph{key} rather
than a value, as in \texttt{\{*-location\}}, asks for a match over label \emph{names},
which a PromQL selector cannot do, because label names are not queryable there ---
Remark~\ref{rem:label-name} puts that to a real server, including the quoted-name
syntax 3.x added, which looks like somewhere a pattern could go and is not.

This paragraph used to argue the fill row as well, and the argument it used is the one
Remark~\ref{rem:fill} broke: that a PromQL expression has no term introducing a sample
where the database holds none, with \texttt{or vector(0)} offered as the nearest thing
and correctly rejected for supplying one unlabelled series rather than a value per
series. The rejection was right and the generalisation was not. What survives of the
paragraph is the observation it made in passing, which turns out to have been the more
useful one: the code's name never fitted its contents. Datadog's interpolating modes
are \texttt{.fill(linear)} and \texttt{.fill(last)}, and this corpus holds one
\texttt{last} and no \texttt{linear} at all --- the other $35$ are constant fills. The
row was named for a construct the corpus does not contain, and a name is a poor thing
to argue from.

\texttt{SCOPE\_DISJUNCTION} is the interesting row, and it is interesting for having
almost emptied. A Datadog scope is a boolean query over tags
(Figure~\ref{fig:productions}) and a PromQL selector is a \emph{conjunction} of matchers,
so disjunction is the one scope operator with no selector-level form --- unless every
branch constrains the same tag key, in which case Proposition~\ref{prop:disjunction}
makes it a single regex matcher. An earlier version of this translator refused all of
it, at $124$ occurrences, and reported that as a gap in the language. It was not: $120$
of those occurrences are single-key disjunctions, and implementing
Proposition~\ref{prop:disjunction} translates them, taking $\kappa$ from $94.2\%$ to
$95.6\%$ at the time, and $96.6\%$ as it now stands.

What survives is four occurrences: disjunction across
different keys (\texttt{\{(status:400 OR s)\}}), a branch that constrains two keys at
once (\texttt{\{(a:1 AND b:1) OR (a:2 AND b:2)\}}), and one De~Morgan case that becomes
the same thing. Those remain sayable in PromQL as a union of selectors, but not as the
syntax-directed translation of a scope --- which is why the split below counts them as
unfinished work rather than as a language gap. That expressibility claim is held to a
weaker standard than the one \texttt{top(...)} is held to, and
the asymmetry should be visible rather than quietly enjoyed: the union there was
written out and read off a real server, while this one is sketched and never built, so
what moves these four occurrences into the engineering column is an argument rather
than a reading. Four occurrences do not move the split, but the standard is the point
of the split. An earlier version of this section put
them in the inexpressible column, where they contradicted the sentence above that says
PromQL expresses them; a property of this translator's architecture is not a property
of the target language.

The episode is worth recording rather than quietly fixing, because the mistake is the
one this paper is most at risk of making throughout. A translator's refusal looks like
evidence about the languages, and it is only evidence about the translator until
something separates the two. Here what separated them was re-parsing the refused queries
and asking which case each was --- a question the gap code alone could not answer, since
one code covered both.

\begin{table}[htbp]\centering\footnotesize
\begin{tabular}{@{}llrr@{}}
\toprule
Feature & Construct & Occurrences & Translated \\
\midrule
\texttt{AGG\_COUNT} & space aggregator count: & 58 & 0.0\% \\
\texttt{FUNC\_APPLIED} & function applied to the query & 205 & 0.0\% \\
\texttt{SCOPE\_BARE} & scope term that is a bare tag key & 92 & 0.0\% \\
\texttt{STR\_SINGLE\_QUOTED} & single-quoted string literal & 139 & 0.0\% \\
\texttt{SCOPE\_PAREN} & parenthesised scope & 122 & 0.0\% \\
\texttt{SCOPE\_IN} & scope membership test IN (...) & 51 & 0.0\% \\
\texttt{M\_WEIGHTED} & .weighted() & 70 & 0.0\% \\
\texttt{AGG\_MAX} & space aggregator max: & 275 & 4.0\% \\
\texttt{SCOPE\_BOOL} & boolean scope (AND / OR / NOT) & 160 & 14.4\% \\
\texttt{AGG\_MIN} & space aggregator min: & 51 & 35.3\% \\
\texttt{CONST} & numeric literal & 470 & 62.8\% \\
\texttt{M\_FILL} & .fill(...) & 185 & 63.2\% \\
\bottomrule
\end{tabular}

\caption{The twelve features with the lowest translated share, among those appearing in
at least $50$ query occurrences. Rows read ``queries using this feature'', not ``this
feature''.}
\label{tab:gap-by-feature}
\end{table}

Table~\ref{tab:gap-by-feature} attributes the refusals to the Datadog features the
refused queries use, and the caveat in its caption is load-bearing. A feature at $0.0\%$
is not necessarily one the translator cannot handle; it can be one that in this corpus
only ever appears beside something else that fails. \texttt{SCOPE\_BARE} demonstrated
this and then stopped: \texttt{avg:x\{host\}} alone translates to
\texttt{x\{host=\textasciitilde{}".+"\}}, yet before
Proposition~\ref{prop:disjunction} the feature stood at $0.0\%$, because in these
dashboards a bare tag key occurs only inside a boolean scope or next to a
\texttt{*-location} key glob. Translating the scopes moved it to $83.6\%$ without anyone
touching bare tag keys. Reading such a row as a statement about its own feature would
have overstated the gaps, which is why the gap catalogue and not this table is the
claim.

